\documentclass{ULBreport}
\usepackage{booktabs}
\usepackage{amsmath}
\usepackage{amssymb}
\usepackage{listings}
\sceau{Pictures/sceauULB.jpg}

\addbibresource{biblio.bib}

\begin{document}

\titleULB{
	title={Assignment 1},
	studies={2026 - 2027},
	course={INFO-H419 - Data Warehouses},
	author={Oliver Irvine\\ Jianli Zhao\\ Sager Singh\\ Neha},
	date={November 2026},
	teacher={Esteban Zimányi},
	logo={Pictures/logo-polytech.jpg},
	manyAuthor}

\chapter{Report}

\section{Generalised Relational and Physical Optimization Framework}

An initial analysis of the 99-query TPC-DS workload across Scale Factors (SF) 1, 2, 4, and 8 revealed that the most severe performance bottlenecks were runtimes jumping by $50\times$ to over $1,900\times$ between SF4 and SF8. They stemm from a recurring set of logical and physical execution flaws in MySQL 8.0. Rather than treating each slow query as an isolated edge case, we identified and formulated three generalisable optimisation rules grounded in relational algebra equivalences and physical storage architecture.

\subsection{Law I: Distributive Aggregate Pushdown (Early Surrogate-Key Grouping)}

\subsubsection{Theoretical Foundation}
Let $R$ be a large fact relation (e.g., \texttt{store\_sales}, \texttt{catalog\_sales}) and $S$ be a dimension relation (e.g., \texttt{customer}, \texttt{item}, \texttt{store}) or a second fact relation (e.g., \texttt{inventory}). Suppose a query computes an aggregate $\gamma_{G,\, \text{AGG}(x)}(R \bowtie_{\theta} S)$, where the grouping set $G$ contains descriptive attributes from $S$ (e.g. customer names) and the join condition $\theta$ only references a subset of join keys $K_R \subseteq \text{attr}(R)$.

By the distributive property of summation and counting over relational joins, we can push the grouping operator $\gamma$ down directly onto $R$ prior to evaluating the join $\bowtie_{\theta}$, provided we group by the compact surrogate key attributes $G_R = (G \cap \text{attr}(R)) \cup K_R$:
\begin{equation}
\gamma_{G,\, \text{SUM}(x),\, \text{COUNT}(*)}\left(R \bowtie_{\theta} S\right) \;\equiv\; \gamma_{G,\, \text{SUM}(\text{sum\_x}),\, \text{SUM}(\text{cnt})}\left( \left(\gamma_{G_R,\, \text{SUM}(x)\to\text{sum\_x},\, \text{COUNT}(*)\to\text{cnt}}(R)\right) \bowtie_{\theta} S \right)
\end{equation}
Furthermore, when $K_R$ is a foreign key referencing the primary key of dimension $S$ (e.g., \texttt{c\_customer\_sk} $\to$ \texttt{c\_customer\_id}, \texttt{c\_first\_name}, \dots), the functional dependency $K_R \to G_S$ guarantees that grouping by the 4-byte integer surrogate key $K_R$ first is 100\% lossless with respect to the outer grouping attributes $G_S$.

\subsubsection{Physical Impact in MySQL}
When a query joins a multi-million-row fact table to wide dimension tables before grouping, MySQL's execution engine must materialise an internal temporary hash aggregation table whose composite hash key consists of multiple variable-length strings (e.g., 7 \texttt{VARCHAR}/\texttt{CHAR} columns in \texttt{customer} or 10 address columns in \texttt{store}). As scale factors grow from SF1 to SF8, this wide temporary table exceeds \texttt{tmp\_table\_size} / \texttt{max\_heap\_table\_size}, spilling to disk and causing severe I/O thrashing. By applying Law I, MySQL aggregates millions of fact rows using compact 4-byte integer surrogate keys (\texttt{\_sk}) in memory, reducing the cardinality of the intermediate relation by one to two orders of magnitude before ever probing the dimension table or reading wide string attributes.

\subsection{Law II: Selection Pushdown and Disjunctive (\texttt{OR}) Bounding-Box Extraction}

\subsubsection{Theoretical Foundation}
Standard relational query optimisers routinely push conjunctive selection predicates ($\sigma_{p_1 \land p_2}(R \bowtie S) \equiv \sigma_{p_1}(R) \bowtie \sigma_{p_2}(S)$) down to base relations. However, in complex analytical queries, selection predicates and equijoin conditions are frequently entangled inside disjunctive (\texttt{OR}) clauses spanning multiple relations:
\begin{equation}
\sigma_{\bigvee_{i=1}^{k} \left(\theta_{\text{join}} \,\land\, p_i(R) \,\land\, q_i(S)\right)}(R \times S)
\end{equation}
By the distributive law of conjunction over disjunction, common equijoin predicates $\theta_{\text{join}}$ can be factored out of the disjunction. More importantly, each relation's disjunctive domain forms a superset ``bounding box'' filter $P_R = \bigvee_{i=1}^{k} p_i(R)$ and $Q_S = \bigvee_{i=1}^{k} q_i(S)$ that must hold for any tuple to survive the overall disjunction:
\begin{equation}
\sigma_{\bigvee_{i=1}^{k} \left(\theta_{\text{join}} \,\land\, p_i(R) \,\land\, q_i(S)\right)}(R \times S) \;\equiv\; \sigma_{\bigvee_{i=1}^{k} \left(p_i(R) \,\land\, q_i(S)\right)}\left( \sigma_{\bigvee_{i=1}^{k} p_i(R)}(R) \;\bowtie_{\theta_{\text{join}}}\; \sigma_{\bigvee_{i=1}^{k} q_i(S)}(S) \right)
\end{equation}
Similarly, when multiple Common Table Expressions (CTEs) scan multiple years of data only for the outer \texttt{WHERE} clause to restrict all CTEs to a single target year, pushing the outer restriction $\sigma_{\text{year}=y}$ directly into each CTE eliminates redundant scans.

\subsubsection{Physical Impact in MySQL}
MySQL's cost-based optimiser does not automatically factor equijoins out of multi-table \texttt{OR} blocks, nor does it synthesise bounding-box pre-filters for the participating dimension and fact tables. Consequently, it evaluates the join across unfiltered base tables (e.g., scanning 23M rows of \texttt{store\_sales} against 1.92M rows of \texttt{customer\_demographics} at SF8) and applies the \texttt{OR} filter post-join. Extracting the bounding box of each table into an isolated CTE restricts dimension tables to a tiny fraction of their rows (e.g., 8.5\% of \texttt{customer\_demographics} in Query 13) and prunes out-of-range fact tuples immediately upon index lookup.

\subsection{Law III: Join Associativity, CTE Materialisation (\texttt{NO\_MERGE}), and Composite PK/FK Driving Order}

\subsubsection{Theoretical Foundation}
By the commutativity and associativity of inner joins, $(F_1 \bowtie F_2) \bowtie \sigma_p(D) \equiv (\sigma_p(D) \bowtie F_1) \bowtie F_2$. While all join trees produce identical relational results, the intermediate cardinalities vary by orders of magnitude depending on whether selective dimension filters $\sigma_p(D)$ are evaluated prior to fact-to-fact joins ($F_1 \bowtie F_2$).

\subsubsection{Physical Impact in MySQL}
Two physical mechanics in MySQL 8.0 InnoDB govern how join trees scale from SF1 to SF8:
\begin{enumerate}
    \item \textbf{Defeating \texttt{derived\_merge} via \texttt{NO\_MERGE()}:} By default, MySQL's \texttt{derived\_merge} optimiser flag dissolves non-aggregated CTEs and derived tables back into the outer query's flat join graph. At SF8, cardinality misestimations across 6--10 tables frequently cause the optimiser to abandon selective dimension filters and initiate a nested-loop join from an unselective table or fact table, creating multi-billion-row intermediate explosions. Adding the \texttt{/*+ NO\_MERGE() */} optimiser hint (or a \texttt{GROUP BY}) inside a filtered CTE forces MySQL to materialise the small filtered result set into an in-memory temporary table first.
    \item \textbf{Driving into InnoDB Clustered Composite Primary Keys and Foreign Key B-Trees:} In InnoDB, every table is physically clustered around its \texttt{PRIMARY KEY}, and every foreign key constraint creates a secondary B-Tree index. In TPC-DS, fact tables have composite primary keys such as \texttt{inventory(inv\_date\_sk, inv\_item\_sk, inv\_warehouse\_sk)} or \texttt{store\_sales(ss\_item\_sk, ss\_ticket\_number)}, as well as foreign key indexes on date surrogate keys (e.g., \texttt{cs\_sold\_date\_sk}, \texttt{ss\_sold\_date\_sk}). By combining materialised selective dimension CTEs with \texttt{STRAIGHT\_JOIN}, we force MySQL to drive the execution plan from the small filtered dimension CTE directly into the fact table's Foreign Key index or left-prefix Composite Primary Key, converting 94-million-row full table scans into $O(\log N)$ B-Tree point lookups.
\end{enumerate}

\section{Group A: Early Surrogate-Key Aggregation Before Dimension and Self-Joins (Law I + Law III)}

Queries in Group A (Q4, Q11, Q21, Q47, Q50, Q74, Q81) aggregate large fact tables (\texttt{store\_sales}, \texttt{catalog\_sales}, \texttt{web\_sales}, \texttt{catalog\_returns}, \texttt{inventory}) by string-valued dimension attributes (\texttt{customer}, \texttt{item}, \texttt{store}, \texttt{warehouse}), or self-join multi-year CTEs on string natural keys. By applying \textbf{Law I} (grouping fact rows by integer surrogate keys before joining dimensions) and \textbf{Law III} (driving from small, materialised \texttt{date\_dim} CTEs into fact-table indexes via \texttt{STRAIGHT\_JOIN}), we remove wide-string temporary tables, redundant self-joins and statistics-dependent plan choices.

\paragraph{Refinement of Law I.} Pushing the aggregation below the dimension join is only lossless when the outer grouping attributes are functionally determined by the surrogate key \emph{and} the query still groups by those attributes afterwards. In the TPC-DS data this cannot be taken for granted: \texttt{c\_customer\_id} is unique per \texttt{c\_customer\_sk}, but \texttt{item}, \texttt{store} and \texttt{warehouse} are slowly changing dimensions. At SF1, 18,000 item keys map to only 9,000 values of \texttt{i\_item\_id}, 12 stores share 8 store names and 5 warehouses share 4 names. Whenever a query groups by such attributes, our rewrites therefore aggregate by surrogate keys first and then \emph{re-aggregate} by the original attributes after the dimension join. This is exact for distributive aggregates: $\text{SUM}(\text{SUM}(x))$ equals $\text{SUM}(x)$, including returning \texttt{NULL} only when all inputs are \texttt{NULL}.

\paragraph{Corrected baseline for Q4 and Q11.} The inherited \texttt{qmod\_combined} versions of Q4 and Q11 did not preserve the original queries: Q4 compared absolute 2001 totals instead of the 2001$\rightarrow$2002 growth ratios of the three channels, and both returned \texttt{customer\_birth\_country} instead of \texttt{customer\_login}. On SF1 the inherited Q4 returns 245 rows, whereas the original returns 8. We replaced both with the MySQL-compatible original queries (which differ from \texttt{orig} only in whitespace); the historical Q4 and Q11 timings refer to the incorrect queries and are not used here.

\paragraph{Experimental setup and correctness.} All Group A measurements were taken on one machine (Apple M4, 10 cores, 24\,GiB RAM, MySQL 8.4.13, 6\,GiB InnoDB buffer pool) on SF1 (4.34\,GB including indexes) and SF2 (7.59\,GB, larger than the buffer pool). Each query was executed once as a warm-up and then timed; we report the median of three such runs, including the time to fetch all rows. Correctness was verified by executing the baseline and the optimised query with the trailing \texttt{LIMIT} removed and comparing the complete result sets (exact comparison of decimals and strings, relative tolerance $10^{-9}$ for floating-point aggregates). All seven queries return identical results, in the same order, on SF1 and SF2.

\subsection{Group A Performance Summary}
Table~\ref{tab:group_a_results} reports the speedup of each optimised query over the baseline and Table~\ref{tab:group_a_times} the underlying execution times. Measurements on SF4 and SF8 are pending.

\begin{table}[htbp]
\centering
\caption{Group A Improvement Factor (Speedup) Over Baseline Across Scale Factors}
\label{tab:group_a_results}
\begin{tabular}{@{}lrrrrr@{}}
\toprule
\textbf{Query} & \textbf{SF1 ($\times$)} & \textbf{SF2 ($\times$)} & \textbf{SF4 ($\times$)} & \textbf{SF8 ($\times$)} & \textbf{Total ($\times$)} \\ \midrule
Query 4  & 16.7 & 19.3 & -- & -- & -- \\
Query 11 & 16.7 & 17.3 & -- & -- & -- \\
Query 21 & 2.4 & 133.1 & -- & -- & -- \\
Query 47 & 5.0 & 5.9 & -- & -- & -- \\
Query 50 & 1.4 & 1.6 & -- & -- & -- \\
Query 74 & 5.8 & 6.5 & -- & -- & -- \\
Query 81 & 20.3 & 40.0 & -- & -- & -- \\ \midrule
\textbf{Group A Total} & \textbf{11.7} & \textbf{14.9} & \textbf{--} & \textbf{--} & \textbf{--} \\ \bottomrule
\end{tabular}
\end{table}

\begin{table}[htbp]
\centering
\caption{Group A Execution Times in Seconds (Median of Three Warm Runs)}
\label{tab:group_a_times}
\begin{tabular}{@{}lrrrr@{}}
\toprule
\textbf{Query} & \textbf{SF1 Baseline} & \textbf{SF1 Optimised} & \textbf{SF2 Baseline} & \textbf{SF2 Optimised} \\ \midrule
Query 4  & 50.954 & 3.053 & 112.460 & 5.821 \\
Query 11 & 35.522 & 2.128 & 71.957 & 4.150 \\
Query 21 & 0.337 & 0.142 & 25.189 & 0.189 \\
Query 47 & 9.365 & 1.884 & 18.915 & 3.219 \\
Query 50 & 0.032 & 0.023 & 0.050 & 0.032 \\
Query 74 & 11.797 & 2.048 & 25.602 & 3.923 \\
Query 81 & 1.981 & 0.098 & 7.012 & 0.175 \\ \midrule
\textbf{Total} & \textbf{110.0} & \textbf{9.4} & \textbf{261.2} & \textbf{17.5} \\ \bottomrule
\end{tabular}
\end{table}

From SF1 to SF2 the summed baseline runtime grows by $2.37\times$, whereas the optimised queries grow by $1.86\times$. The group totals are dominated by Q4 and Q11, while the largest single improvement is Q21, discussed next.

\subsection{Group A Highlight: Query 21 Optimization}

\subsubsection{Baseline Execution Plan and Bottleneck Analysis}
Query 21 compares, per warehouse and item, the inventory on hand in the 30 days before and after 2002-03-20 for items priced between 0.99 and 1.49. In the baseline (\texttt{qmod\_combined/query21.sql}) it takes only 0.34\,s on SF1 but 25.19\,s on SF2: a $74.7\times$ slowdown for twice the data. The historical results of the full workload show the same jump (0.30\,s on SF1, 53.31\,s on SF2). The cause is a plan flip between the two scale factors.

On SF1 (Listing~\ref{lst:q21_explain_sf1}) MySQL drives from the 61 qualifying \texttt{date\_dim} rows into the left prefix \texttt{inv\_date\_sk} of the \texttt{inventory} primary key \texttt{(inv\_date\_sk, inv\_item\_sk, inv\_warehouse\_sk)}, so only the weekly inventory snapshots inside the window are read.

\begin{lstlisting}[caption={Baseline \texttt{EXPLAIN FORMAT=TREE} on SF1 for Query 21 (top of plan omitted)}, label={lst:q21_explain_sf1}, basicstyle=\ttfamily\scriptsize, breaklines=true, frame=single]
-> Aggregate using temporary table
    -> Inner hash join (warehouse.w_warehouse_sk = inventory.inv_warehouse_sk)  (cost=1.18e+6 rows=255824)
        -> Table scan on warehouse  (cost=397e-6 rows=5)
        -> Hash
            -> Nested loop inner join  (cost=1.05e+6 rows=255824)
                -> Nested loop inner join  (cost=248351 rows=2.3e+6)
                    -> Filter: (date_dim.d_date between <cache>((cast('2002-03-20' as date) - interval 30 day)) and <cache>((cast('2002-03-20' as date) + interval 30 day)))  (cost=7429 rows=8051)
                        -> Table scan on date_dim  (cost=7429 rows=72468)
                    -> Index lookup on inventory using PRIMARY (inv_date_sk=date_dim.d_date_sk)  (cost=1.33 rows=286)
                -> Filter: (item.i_current_price between 0.99 and 1.49)  (cost=0.25 rows=0.111)
                    -> Single-row index lookup on item using PRIMARY (i_item_sk=inventory.inv_item_sk)  (cost=0.25 rows=1)
\end{lstlisting}

On SF2 (Listing~\ref{lst:q21_explain_sf2}) the optimiser instead starts from \texttt{warehouse} and reads \texttt{inventory} through the foreign-key index \texttt{inv\_w}, checking the date of each row only at the end of the join. The reason is a cardinality misestimate: InnoDB's sampled index statistics report 2,925 distinct values for \texttt{inv\_w} on SF2 (3 on SF1), although \texttt{inventory} references only 5 warehouses. MySQL therefore expects about 5,790 rows per warehouse, while each warehouse actually has 3,393,261. The plan reads all 16,966,305 inventory rows across all 261 weekly snapshots, although only 9 snapshots (585,045 rows, 3.4\%) fall inside the date window.

\begin{lstlisting}[caption={Baseline \texttt{EXPLAIN FORMAT=TREE} on SF2 for Query 21 (top of plan omitted)}, label={lst:q21_explain_sf2}, basicstyle=\ttfamily\scriptsize, breaklines=true, frame=single]
-> Aggregate using temporary table
    -> Nested loop inner join  (cost=21395 rows=357)
        -> Nested loop inner join  (cost=20269 rows=3217)
            -> Nested loop inner join  (cost=10136 rows=28952)
                -> Table scan on warehouse  (cost=0.75 rows=5)
                -> Index lookup on inventory using inv_w (inv_warehouse_sk=warehouse.w_warehouse_sk)  (cost=1564 rows=5790)
            -> Filter: (item.i_current_price between 0.99 and 1.49)  (cost=0.25 rows=0.111)
                -> Single-row index lookup on item using PRIMARY (i_item_sk=inventory.inv_item_sk)  (cost=0.25 rows=1)
        -> Filter: (date_dim.d_date between <cache>((cast('2002-03-20' as date) - interval 30 day)) and <cache>((cast('2002-03-20' as date) + interval 30 day)))  (cost=0.25 rows=0.111)
            -> Single-row index lookup on date_dim using PRIMARY (d_date_sk=inventory.inv_date_sk)  (cost=0.25 rows=1)
\end{lstlisting}

\subsubsection{Precise Application of Optimization Laws in SQL}
We applied \textbf{Law III} and \textbf{Law I} to Query 21 as follows:
\begin{enumerate}
    \item \textbf{Application of Law III (Materialised Date Driving Order):} The 61-day window is materialised as a CTE and joined to \texttt{inventory} with \texttt{STRAIGHT\_JOIN}, which fixes the date-driven access path through the primary key prefix regardless of the index statistics. The price filter becomes a semi-join on a small materialised set of item keys.

    \textit{Before (Baseline):}
\begin{lstlisting}[language=SQL, basicstyle=\ttfamily\scriptsize, breaklines=true, frame=single]
FROM inventory, warehouse, item, date_dim
WHERE i_current_price BETWEEN 0.99 AND 1.49
  AND i_item_sk = inv_item_sk
  AND inv_warehouse_sk = w_warehouse_sk
  AND inv_date_sk = d_date_sk
  AND d_date BETWEEN (CAST('2002-03-20' AS DATE) - INTERVAL 30 DAY)
                 AND (CAST('2002-03-20' AS DATE) + INTERVAL 30 DAY)
GROUP BY w_warehouse_name, i_item_id
\end{lstlisting}

    \textit{After (Law III Applied):}
\begin{lstlisting}[language=SQL, basicstyle=\ttfamily\scriptsize, breaklines=true, frame=single]
WITH d_window AS (
    SELECT /*+ NO_MERGE() */ d_date_sk, d_date
    FROM date_dim
    WHERE d_date BETWEEN (CAST('2002-03-20' AS DATE) - INTERVAL 30 DAY)
                     AND (CAST('2002-03-20' AS DATE) + INTERVAL 30 DAY)
),
priced_items AS (
    SELECT /*+ NO_MERGE() */ i_item_sk, i_item_id
    FROM item WHERE i_current_price BETWEEN 0.99 AND 1.49
), ...
    FROM d_window d
    STRAIGHT_JOIN inventory inv ON inv.inv_date_sk = d.d_date_sk
    WHERE inv.inv_item_sk IN (SELECT i_item_sk FROM priced_items)
\end{lstlisting}

    \item \textbf{Application of Law I (Surrogate-Key Aggregation with Re-Aggregation):} The before/after sums are computed per integer key pair \texttt{(inv\_warehouse\_sk, inv\_item\_sk)}. Only then are \texttt{warehouse} and the item names joined, and because neither \texttt{w\_warehouse\_name} nor \texttt{i\_item\_id} is unique per key, the partial sums are aggregated again by these attributes before the ratio filter is applied.

    \textit{After (Law I Applied):}
\begin{lstlisting}[language=SQL, basicstyle=\ttfamily\scriptsize, breaklines=true, frame=single]
inv_totals AS (
    SELECT /*+ NO_MERGE() */ inv.inv_warehouse_sk, inv.inv_item_sk,
           SUM(CASE WHEN CAST(d.d_date AS DATE) <  CAST('2002-03-20' AS DATE)
                    THEN inv.inv_quantity_on_hand ELSE 0 END) AS inv_before,
           SUM(CASE WHEN CAST(d.d_date AS DATE) >= CAST('2002-03-20' AS DATE)
                    THEN inv.inv_quantity_on_hand ELSE 0 END) AS inv_after
    ...
    GROUP BY inv.inv_warehouse_sk, inv.inv_item_sk
)
SELECT * FROM (
    SELECT w.w_warehouse_name, i.i_item_id,
           SUM(t.inv_before) AS inv_before, SUM(t.inv_after) AS inv_after
    FROM inv_totals t
    STRAIGHT_JOIN warehouse w ON w.w_warehouse_sk = t.inv_warehouse_sk
    STRAIGHT_JOIN priced_items i ON i.i_item_sk = t.inv_item_sk
    GROUP BY w.w_warehouse_name, i.i_item_id
) x
WHERE (CASE WHEN inv_before > 0 THEN inv_after / inv_before ELSE NULL END)
      BETWEEN 2.0 / 3.0 AND 3.0 / 2.0
\end{lstlisting}
\end{enumerate}

Listing~\ref{lst:q21_explain_opt} shows the optimised plan on SF2: the date window drives into the \texttt{inventory} primary key, and \texttt{warehouse} and \texttt{item} are only probed once per aggregated key pair. The runtime on SF2 drops from 25.19\,s to 0.19\,s ($133.1\times$), and the optimised query grows only from 0.14\,s to 0.19\,s between SF1 and SF2.

\begin{lstlisting}[caption={Optimised \texttt{EXPLAIN FORMAT=TREE} on SF2 for Query 21 (top of plan omitted)}, label={lst:q21_explain_opt}, basicstyle=\ttfamily\scriptsize, breaklines=true, frame=single]
-> Aggregate using temporary table
    -> Nested loop inner join  (cost=137526 rows=0)
        -> Nested loop inner join  (cost=68764 rows=0)
            -> Table scan on t  (cost=2.5..2.5 rows=0)
                -> Materialize CTE inv_totals  (cost=0..0 rows=0)
                    -> Aggregate using temporary table
                        -> Nested loop inner join  (cost=1.13e+6 rows=275055)
                            -> Nested loop inner join  (cost=258668 rows=2.48e+6)
                                -> Table scan on d  (cost=8220..8323 rows=8051)
                                    -> Materialize CTE d_window  (cost=8220..8220 rows=8051)
                                -> Index lookup on inv using PRIMARY (inv_date_sk=d.d_date_sk)  (cost=0.348 rows=308)
                            -> Filter: (item.i_current_price between 0.99 and 1.49)  (cost=0.25 rows=0.111)
                                -> Single-row index lookup on item using PRIMARY (i_item_sk=inv.inv_item_sk)  (cost=0.25 rows=1)
            -> Single-row index lookup on w using PRIMARY (w_warehouse_sk=t.inv_warehouse_sk)  (cost=0.25 rows=1)
        -> Filter: (item.i_current_price between 0.99 and 1.49)  (cost=0.25 rows=0.111)
            -> Single-row index lookup on item using PRIMARY (i_item_sk=t.inv_item_sk)  (cost=0.25 rows=1)
\end{lstlisting}

\subsection{Other Group A Queries}

\paragraph{Q4, Q11 and Q74: single-pass year pivot.} The baseline materialises a CTE that aggregates every sales channel by seven customer string columns and year (over all years for Q4 and Q11) and then self-joins it four times (six for Q4) on the string \texttt{customer\_id}. The rewrite scans each channel once, driven from the two target years of \texttt{date\_dim} through the sold-date foreign-key index (Law III), and groups by the integer customer key, computing both years with conditional aggregates such as \texttt{SUM(CASE WHEN d\_year = 2001 THEN ... END)} (Law I). The channels are then joined on \texttt{customer\_sk}, and \texttt{customer} is joined once for the output columns. Omitting the \texttt{ELSE} branch keeps a year whose values are all \texttt{NULL} at \texttt{NULL}, and a per-year \texttt{COUNT} preserves the requirement of the original self-joins that the customer has rows in both years. \texttt{STDDEV\_SAMP} in Q74 cannot be merged from partial results, but each (customer, year) group is still computed in a single aggregation. The speedups on SF2 are $19.3\times$ (Q4), $17.3\times$ (Q11) and $6.5\times$ (Q74).

\paragraph{Q47: integer group identifiers and window functions.} A first attempt that pre-aggregated \texttt{store\_sales} by \texttt{(ss\_item\_sk, ss\_store\_sk, month)} gained only about $1.5\times$: on SF2 this key barely reduced the rows (1.34M to 0.96M), and the second aggregation by the string attributes added another full pass. Law I only pays off when the surrogate-key grouping is clearly coarser than the original grouping. The final version maps \texttt{(i\_category, i\_brand)} and \texttt{(s\_store\_name, s\_company\_name)} to integer group identifiers with \texttt{DENSE\_RANK} (which, like \texttt{GROUP BY}, compares strings by collation), aggregates \texttt{store\_sales} once by these integers directly into 76k groups, and replaces the three-way self-join for the previous and next month with \texttt{LAG}/\texttt{LEAD}. Since the original self-join uses equality predicates that never match \texttt{NULL}, partitions with a \texttt{NULL} category, brand, store or company name are excluded in the rewrite as well. The speedup is $5.0\times$ on SF1 and $5.9\times$ on SF2.

\paragraph{Q50: fixed return-driven plan.} Q50 is already fast on SF1 and SF2 (0.03--0.05\,s); its historical degradation (0.15\,s on SF4, 75.81\,s on SF8) only appears at SF8. The rewrite fixes the access path from the 31 days of August 2001 into \texttt{store\_returns} through its returned-date index and into \texttt{store\_sales} through its primary key \texttt{(ss\_item\_sk, ss\_ticket\_number)}, counts the return intervals per \texttt{ss\_store\_sk} and only then re-aggregates by the ten store attributes. Its benefit is expected at SF8 and remains to be measured.

\paragraph{Q81: decorrelated per-state threshold.} The baseline evaluates the average return per state as a dependent subquery, rescanning the matching CTE rows for every outer row, so its cost grows roughly quadratically with the data (17.72\,s on SF4, 507.15\,s on SF8 in the historical results). The rewrite materialises the per-customer returns once, driven from the 2002 dates through the \texttt{catalog\_returns} returned-date index, pre-aggregates the threshold $1.2 \times \text{AVG}$ per state (Law I) and joins it on the state. A \texttt{NULL} state matches neither the original predicate nor the join. The outer filter on \texttt{ca\_state = 'NH'} cannot be pushed into the CTE, because it refers to the customer's current address and not to the address of the return. The speedup is $20.3\times$ on SF1 and $40.0\times$ on SF2.

\section{Group B: Fact-to-Fact Multi-Way Joins and Composite Primary Key Lookups (Law I + Law III)}

Queries in Group B join two to six fact tables together (e.g., \texttt{catalog\_sales} with \texttt{inventory}, or \texttt{store\_sales} with \texttt{store\_returns} and \texttt{catalog\_sales}) across multiple \texttt{date\_dim} aliases. Without intervention, MySQL joins the massive fact tables before filtering by dates or fails to utilise the left-prefix of InnoDB's composite primary keys, causing severe intermediate row explosions.

\subsection{Group B Performance Summary}
Table~\ref{tab:group_b_results} reports the speedup factors achieved over the baseline across all four scale factors for Group B queries.

\begin{table}[htbp]
\centering
\caption{Group B Improvement Factor (Speedup) Over Baseline Across Scale Factors}
\label{tab:group_b_results}
\begin{tabular}{@{}lrrrrr@{}}
\toprule
\textbf{Query} & \textbf{SF1 ($\times$)} & \textbf{SF2 ($\times$)} & \textbf{SF4 ($\times$)} & \textbf{SF8 ($\times$)} & \textbf{Total ($\times$)} \\ \midrule
Query 17 & -- & -- & -- & -- & -- \\
Query 25 & -- & -- & -- & -- & -- \\
Query 29 & -- & -- & -- & -- & -- \\
Query 72 & $21.91\times$ & $92.98\times$ & $80.82\times$ & $142.30\times$ & -- \\
Query 77 & -- & -- & -- & -- & -- \\
Query 80 & -- & -- & -- & -- & -- \\ \midrule
\textbf{Group B Total / Average} & \textbf{--} & \textbf{--} & \textbf{--} & \textbf{--} & \textbf{--} \\ \bottomrule
\end{tabular}
\end{table}

\subsection{Group B Highlight: Query 72 Optimization}

\subsubsection{Baseline Execution Plan and Bottleneck Analysis}
In our baseline Power Test across Scale Factors 1, 2, 4, and 8, \textbf{Query 72} was the single largest performance bottleneck in the entire 99-query workload, requiring \textbf{110.44s} at SF1 and degrading by $48.93\times$ to \textbf{5,404.38s (1.5 hours)} at SF8. In the baseline query, MySQL's \texttt{derived\_merge} optimiser dissolved the CTE boundaries and initiated a nested-loop join from \texttt{warehouse} and \texttt{date\_dim d2} into \texttt{inventory} (94M rows at SF8) and then probed \texttt{catalog\_sales} by \texttt{cs\_item\_sk} alone, generating \textbf{2.17 billion intermediate rows} before checking whether \texttt{cs\_sold\_date\_sk} was in 2002.

Listing~\ref{lst:q72_explain_bad} shows the \texttt{EXPLAIN FORMAT=TREE} on SF8 when CTE merging and nested-loop explosion occur.

\begin{lstlisting}[caption={Unmaterialised \texttt{EXPLAIN FORMAT=TREE} on SF8 for Query 72 Showing 2.17B-Row Explosion}, label={lst:q72_explain_bad}, basicstyle=\ttfamily\scriptsize, breaklines=true, frame=single]
-> Limit: 100 row(s)
    -> Table scan on <temporary>
        -> Aggregate using temporary table
            -> Nested loop inner join  (cost=42.8e+9 rows=36.2e+6)
                -> Nested loop inner join  (cost=42.8e+9 rows=109e+6)
                    -> Nested loop left join  (cost=42.6e+9 rows=2.17e+9)
                        -> Nested loop inner join  (cost=42.4e+9 rows=2.17e+9)
                            -> Nested loop inner join  (cost=41.8e+9 rows=79.3e+6)
                                -> Nested loop inner join  (cost=41.7e+9 rows=79.3e+6)
                                    -> Inner hash join (no condition)  (cost=94.5e+6 rows=420e+6)
                                        -> Table scan on w  (cost=340e-6 rows=8)
                                        -> Hash
                                            -> Inner hash join (d2.d_week_seq = `<subquery6>`.d_week_seq)  (cost=52.5e+6 rows=52.5e+6)
                                                -> Table scan on d2  (cost=2993 rows=72468)
                                                -> Hash
                                                    -> Table scan on <subquery6>  (cost=9583..9676 rows=7247)
                                    -> Filter: (inv.inv_warehouse_sk = w.w_warehouse_sk)  (cost=2.21 rows=0.189)
                                        -> Index lookup on inv using PRIMARY (inv_date_sk = d2.d_date_sk)  (cost=2.21 rows=991)
                                -> Single-row index lookup on i using PRIMARY (i_item_sk = inv.inv_item_sk)  (cost=0.983 rows=1)
                            -> Filter: ((inv.inv_quantity_on_hand < cs.cs_quantity) and (cs.cs_sold_date_sk is not null) and (cs.cs_ship_date_sk is not null))  (cost=1.44 rows=27.4)
                                -> Index lookup on cs using PRIMARY (cs_item_sk = inv.inv_item_sk)  (cost=1.44 rows=82.1)
                        -> Single-row covering index lookup on p using PRIMARY (p_promo_sk = cs.cs_promo_sk)  (cost=0.893 rows=1)
                    -> Filter: ((date_dim.d_year = 2002) and (date_dim.d_week_seq = `<subquery6>`.d_week_seq))  (cost=0.99 rows=0.05)
                        -> Single-row index lookup on date_dim using PRIMARY (d_date_sk = cs.cs_sold_date_sk)  (cost=0.99 rows=1)
                -> Filter: (cast(d3.d_date as double) > cast((date_dim.d_date + 5) as double))  (cost=0.99 rows=0.333)
                    -> Single-row index lookup on d3 using PRIMARY (d_date_sk = cs.cs_ship_date_sk)  (cost=0.99 rows=1)
\end{lstlisting}

\subsubsection{Precise Application of Optimization Laws in SQL}
We applied \textbf{Law I} and \textbf{Law III} to Query 72 as follows:
\begin{enumerate}
    \item \textbf{Application of Law I (Distributive Pre-Aggregation of \texttt{catalog\_sales}):} Any two \texttt{catalog\_sales} rows in the same week that share $(\texttt{cs\_item\_sk}, \texttt{d\_week\_seq}, \texttt{cs\_quantity})$ match the exact same set of \texttt{inventory} records under $(\texttt{cs\_item\_sk} = \texttt{inv\_item\_sk} \land \texttt{d1.d\_week\_seq} = \texttt{d2.d\_week\_seq} \land \texttt{inv\_quantity\_on\_hand} < \texttt{cs\_quantity})$. Pre-grouping \texttt{filtered\_cs} by these three integer columns collapses duplicate item-week-quantity sales before probing \texttt{inventory}, and delays joining the wide \texttt{item} (\texttt{VARCHAR(200)}) and \texttt{warehouse} strings until after the fact join.

    \textit{Before (Baseline):}
\begin{lstlisting}[language=SQL, basicstyle=\ttfamily\scriptsize, breaklines=true, frame=single]
SELECT i_item_desc, w_warehouse_name, d1.d_week_seq,
       SUM(CASE WHEN p_promo_sk IS NULL THEN 1 ELSE 0 END) no_promo,
       SUM(CASE WHEN p_promo_sk IS NOT NULL THEN 1 ELSE 0 END) promo,
       COUNT(*) total_cnt
FROM catalog_sales
JOIN inventory ON (cs_item_sk = inv_item_sk AND inv_quantity_on_hand < cs_quantity)
JOIN warehouse ON (w_warehouse_sk = inv_warehouse_sk)
JOIN item ON (i_item_sk = cs_item_sk)
...
LEFT OUTER JOIN promotion ON (cs_promo_sk = p_promo_sk)
\end{lstlisting}

    \textit{After (Law I Applied):}
\begin{lstlisting}[language=SQL, basicstyle=\ttfamily\scriptsize, breaklines=true, frame=single]
filtered_cs AS (
    SELECT /*+ NO_MERGE() */
           cs.cs_item_sk, d1.d_week_seq, cs.cs_quantity,
           SUM(CASE WHEN p.p_promo_sk IS NULL THEN 1 ELSE 0 END) AS no_promo_cnt,
           SUM(CASE WHEN p.p_promo_sk IS NOT NULL THEN 1 ELSE 0 END) AS promo_cnt,
           COUNT(*) AS cs_cnt
    FROM d1_2002 d1
    STRAIGHT_JOIN catalog_sales cs ON cs.cs_sold_date_sk = d1.d_date_sk
    STRAIGHT_JOIN date_dim d3 ON cs.cs_ship_date_sk = d3.d_date_sk AND d3.d_date > d1.d_date + 5
    LEFT OUTER JOIN promotion p ON cs.cs_promo_sk = p.p_promo_sk
    GROUP BY cs.cs_item_sk, d1.d_week_seq, cs.cs_quantity
)
\end{lstlisting}

    \item \textbf{Application of Law III (Materialised CTEs \& 2-Column Composite PK Lookup into \texttt{inventory}):} By enforcing \texttt{NO\_MERGE()} on \texttt{d1\_2002} and \texttt{filtered\_cs} and joining \texttt{filtered\_cs} to \texttt{d2\_2002} before probing \texttt{inventory}, MySQL supplies both \texttt{d2.d\_date\_sk} and \texttt{cs.cs\_item\_sk} simultaneously into \texttt{inventory}'s composite primary key \texttt{PRIMARY KEY (inv\_date\_sk, inv\_item\_sk, inv\_warehouse\_sk)}, reducing each probe to a direct B-Tree lookup averaging 8.08 rows instead of a 94M-row scan.

    \textit{Before (Baseline):}
\begin{lstlisting}[language=SQL, basicstyle=\ttfamily\scriptsize, breaklines=true, frame=single]
JOIN date_dim d1 ON (cs_sold_date_sk = d1.d_date_sk)
JOIN date_dim d2 ON (inv_date_sk = d2.d_date_sk AND d1.d_week_seq = d2.d_week_seq)
WHERE d1.d_year = 2002
\end{lstlisting}

    \textit{After (Law III Applied):}
\begin{lstlisting}[language=SQL, basicstyle=\ttfamily\scriptsize, breaklines=true, frame=single]
FROM filtered_cs cs
STRAIGHT_JOIN filtered_inv inv
  ON cs.cs_item_sk = inv.inv_item_sk
 AND cs.d_week_seq = inv.d_week_seq
 AND inv.inv_quantity_on_hand < cs.cs_quantity
STRAIGHT_JOIN warehouse w ON w.w_warehouse_sk = inv.inv_warehouse_sk
STRAIGHT_JOIN item i ON i.i_item_sk = cs.cs_item_sk
\end{lstlisting}
\end{enumerate}

Listing~\ref{lst:q72_explain_good} shows the resulting \texttt{EXPLAIN FORMAT=TREE} on SF8, confirming that \texttt{filtered\_cs} is materialised via foreign-key index \texttt{cs\_d2} and that \texttt{inventory} is probed via its 2-column composite primary key (\texttt{inv\_date\_sk = d2.d\_date\_sk, inv\_item\_sk = cs.cs\_item\_sk}, \texttt{rows=8.08}), reducing SF8 execution time from \textbf{5,404.38s to 37.98s} ($142.30\times$ speedup).

\begin{lstlisting}[caption={Optimized \texttt{EXPLAIN FORMAT=TREE} on SF8 for Query 72}, label={lst:q72_explain_good}, basicstyle=\ttfamily\scriptsize, breaklines=true, frame=single]
-> Limit: 100 row(s)
    -> Sort: total_cnt DESC, i.i_item_desc, w.w_warehouse_name, cs.d_week_seq, limit input to 100 row(s) per chunk
        -> Table scan on <temporary>
            -> Aggregate using temporary table
                -> Nested loop inner join
                    -> Nested loop inner join
                        -> Nested loop inner join
                            -> Nested loop inner join
                                -> Inner hash join (no condition)
                                    -> Filter: (cs.d_week_seq is not null)
                                        -> Table scan on cs
                                            -> Materialize CTE filtered_cs
                                                -> Table scan on <temporary>
                                                    -> Aggregate using temporary table
                                                        -> Nested loop inner join
                                                            -> Nested loop left join
                                                                -> Nested loop inner join
                                                                    -> Table scan on d1 (Materialize CTE d1_2002)
                                                                    -> Index lookup on cs using cs_d2 (cs_sold_date_sk = d1.d_date_sk)
                                                                -> Single-row covering index lookup on p using PRIMARY (p_promo_sk = cs.cs_promo_sk)
                                                            -> Filter: (cast(d3.d_date as double) > cast((d1.d_date + 5) as double))
                                                                -> Single-row index lookup on d3 using PRIMARY (d_date_sk = cs.cs_ship_date_sk)
                                    -> Hash (Constant row from EXISTS subqueries)
                                -> Index lookup on d2 using <auto_key0> (d_week_seq = cs.d_week_seq) (Materialize CTE d2_2002)
                            -> Filter: (inv.inv_quantity_on_hand < cs.cs_quantity)  (cost=1.01 rows=2.69)
                                -> Index lookup on inv using PRIMARY (inv_date_sk = d2.d_date_sk, inv_item_sk = cs.cs_item_sk)  (cost=1.01 rows=8.08)
                        -> Single-row index lookup on w using PRIMARY (w_warehouse_sk = inv.inv_warehouse_sk)
                    -> Single-row index lookup on i using PRIMARY (i_item_sk = cs.cs_item_sk)
\end{lstlisting}

\section{Group C: Disjunctive (\texttt{OR}) Bounding-Box Extraction and Star-Join Ordering (Law II + Law III)}

Queries in Group C exhibit severe plan flips between SF4 and SF8 (degrading by $16\times$ to $2,773\times$) because equijoin predicates and dimension filters are trapped inside multi-branch disjunctive (\texttt{OR}) clauses, or because MySQL abandons selective dimension indexes in multi-table star joins. Applying \textbf{Law II} (factoring out equijoins and extracting bounding-box predicates into CTEs) together with \textbf{Law III} (materializing filtered dimensions with \texttt{NO\_MERGE()} and driving into fact indexes via \texttt{STRAIGHT\_JOIN}) restores linear scaling.

\subsection{Group C Performance Summary}
Table~\ref{tab:group_c_results} provides the speedup table for Group C queries across SF1, SF2, SF4, SF8, and Total.

\begin{table}[htbp]
\centering
\caption{Group C Improvement Factor (Speedup) Over Baseline Across Scale Factors}
\label{tab:group_c_results}
\begin{tabular}{@{}lrrrrr@{}}
\toprule
\textbf{Query} & \textbf{SF1 ($\times$)} & \textbf{SF2 ($\times$)} & \textbf{SF4 ($\times$)} & \textbf{SF8 ($\times$)} & \textbf{Total ($\times$)} \\ \midrule
Query 7  & -- & -- & -- & -- & -- \\
Query 13 & -- & -- & -- & -- & -- \\
Query 19 & -- & -- & -- & -- & -- \\
Query 27 & -- & -- & -- & -- & -- \\
Query 36 & -- & -- & -- & -- & -- \\
Query 43 & -- & -- & -- & -- & -- \\
Query 46 & -- & -- & -- & -- & -- \\
Query 48 & -- & -- & -- & -- & -- \\
Query 53 & -- & -- & -- & -- & -- \\
Query 63 & -- & -- & -- & -- & -- \\
Query 73 & -- & -- & -- & -- & -- \\
Query 79 & -- & -- & -- & -- & -- \\
Query 89 & -- & -- & -- & -- & -- \\ \midrule
\textbf{Group C Total / Average} & \textbf{--} & \textbf{--} & \textbf{--} & \textbf{--} & \textbf{--} \\ \bottomrule
\end{tabular}
\end{table}

\subsection{Group C Highlight: Query 13 Optimization}

\subsubsection{Baseline Execution Plan and Bottleneck Analysis}
In the baseline benchmark, \textbf{Query 13} was the second slowest query at SF8, jumping from 3.67s at SF1 and 17.96s at SF4 to \textbf{715.51s} at SF8 ($194.98\times$ degradation). Inspection of \texttt{qmod\_combined/query13.sql} reveals that the equijoin conditions \texttt{ss\_hdemo\_sk = hd\_demo\_sk}, \texttt{cd\_demo\_sk = ss\_cdemo\_sk}, and \texttt{ss\_addr\_sk = ca\_address\_sk} are repeated inside each branch of two 3-way \texttt{OR} blocks. Consequently, MySQL fails to push the dimension filters down before joining, scanning millions of rows across \texttt{store\_sales} ($\approx 23\text{M}$ rows at SF8), \texttt{customer\_demographics} ($1.92\text{M}$ rows), and \texttt{customer\_address} ($400\text{K}$ rows).

Listing~\ref{lst:q13_explain_orig} shows the baseline \texttt{EXPLAIN FORMAT=TREE} on SF8.

\begin{lstlisting}[caption={Baseline \texttt{EXPLAIN FORMAT=TREE} on SF8 for Query 13}, label={lst:q13_explain_orig}, basicstyle=\ttfamily\scriptsize, breaklines=true, frame=single]
% TODO: Paste Baseline Query 13 EXPLAIN FORMAT=TREE output on SF8 here
\end{lstlisting}

\subsubsection{Precise Application of Optimization Laws in SQL}
We applied \textbf{Law II} and \textbf{Law III} to Query 13 as follows:
\begin{enumerate}
    \item \textbf{Application of Law II (Equijoin Factoring \& Disjunctive Bounding-Box Extraction):} Across the three demographic branches, \texttt{hd\_dep\_count} only ever takes values in $\{1, 3\}$ (20\% of \texttt{household\_demographics}), \texttt{(cd\_marital\_status, cd\_education\_status)} only ever matches 3 out of 35 possible combinations (8.5\% of \texttt{customer\_demographics}), and \texttt{ss\_sales\_price} is bounded by $[50.00, 200.00]$. Similarly, across the three address branches, \texttt{ca\_country = 'United States'}, \texttt{ca\_state} is bounded by the 9 target states, and \texttt{ss\_net\_profit} is bounded by $[50, 300]$. We factor out all equijoins and push the bounding-box domain filters directly into the base relations.

    \textit{Before (Baseline):}
\begin{lstlisting}[language=SQL, basicstyle=\ttfamily\scriptsize, breaklines=true, frame=single]
AND ((ss_hdemo_sk = hd_demo_sk AND cd_demo_sk = ss_cdemo_sk
      AND cd_marital_status = 'M' AND cd_education_status = 'Secondary'
      AND ss_sales_price BETWEEN 100.00 AND 150.00 AND hd_dep_count = 3)
  OR (ss_hdemo_sk = hd_demo_sk AND cd_demo_sk = ss_cdemo_sk
      AND cd_marital_status = 'U' AND cd_education_status = '4 yr Degree'
      AND ss_sales_price BETWEEN 50.00 AND 100.00 AND hd_dep_count = 1)
  OR (ss_hdemo_sk = hd_demo_sk AND cd_demo_sk = ss_cdemo_sk
      AND cd_marital_status = 'W' AND cd_education_status = '2 yr Degree'
      AND ss_sales_price BETWEEN 150.00 AND 200.00 AND hd_dep_count = 1))
\end{lstlisting}

    \textit{After (Law II Applied):}
\begin{lstlisting}[language=SQL, basicstyle=\ttfamily\scriptsize, breaklines=true, frame=single]
filtered_hd AS (
    SELECT /*+ NO_MERGE() */ hd_demo_sk, hd_dep_count
    FROM household_demographics WHERE hd_dep_count IN (1, 3)
),
filtered_cd AS (
    SELECT /*+ NO_MERGE() */ cd_demo_sk, cd_marital_status, cd_education_status
    FROM customer_demographics
    WHERE (cd_marital_status = 'M' AND cd_education_status = 'Secondary')
       OR (cd_marital_status = 'U' AND cd_education_status = '4 yr Degree')
       OR (cd_marital_status = 'W' AND cd_education_status = '2 yr Degree')
),
filtered_ca AS (
    SELECT /*+ NO_MERGE() */ ca_address_sk, ca_state
    FROM customer_address
    WHERE ca_country = 'United States'
      AND ca_state IN ('IA', 'MO', 'TX', 'GA', 'LA', 'SD', 'NY', 'TN', 'VA')
)
...
WHERE ss.ss_sales_price BETWEEN 50.00 AND 200.00
  AND ss.ss_net_profit BETWEEN 50 AND 300
  AND ( ... /* exact branch correlation checks */ )
\end{lstlisting}

    \item \textbf{Application of Law III (Materialised Dimension Driving Order):} By materializing \texttt{filtered\_dates} (\texttt{d\_year = 2001}) with \texttt{NO\_MERGE()} and using \texttt{STRAIGHT\_JOIN}, MySQL drives from the 365 days of 2001 into \texttt{store\_sales} via its sold-date foreign key index, filters \texttt{ss\_sales\_price} and \texttt{ss\_net\_profit} immediately on the fact row, and probes the tiny materialised dimension tables in memory.

    \textit{Before (Baseline):}
\begin{lstlisting}[language=SQL, basicstyle=\ttfamily\scriptsize, breaklines=true, frame=single]
FROM store_sales, store, customer_demographics,
     household_demographics, customer_address, date_dim
WHERE s_store_sk = ss_store_sk
  AND ss_sold_date_sk = d_date_sk AND d_year = 2001
\end{lstlisting}

    \textit{After (Law III Applied):}
\begin{lstlisting}[language=SQL, basicstyle=\ttfamily\scriptsize, breaklines=true, frame=single]
FROM filtered_dates d
STRAIGHT_JOIN store_sales ss ON ss.ss_sold_date_sk = d.d_date_sk
STRAIGHT_JOIN store s ON s.s_store_sk = ss.ss_store_sk
STRAIGHT_JOIN filtered_hd hd ON ss.ss_hdemo_sk = hd.hd_demo_sk
STRAIGHT_JOIN filtered_cd cd ON ss.ss_cdemo_sk = cd.cd_demo_sk
STRAIGHT_JOIN filtered_ca ca ON ss.ss_addr_sk = ca.ca_address_sk
\end{lstlisting}
\end{enumerate}

Listing~\ref{lst:q13_explain_opt} shows the resulting \texttt{EXPLAIN FORMAT=TREE} on SF8 for the optimised Query 13.

\begin{lstlisting}[caption={Optimised \texttt{EXPLAIN FORMAT=TREE} on SF8 for Query 13}, label={lst:q13_explain_opt}, basicstyle=\ttfamily\scriptsize, breaklines=true, frame=single]
% TODO: Paste Optimised Query 13 EXPLAIN FORMAT=TREE output on SF8 here
\end{lstlisting}

\end{document}